\documentclass[10.5pt]{article}

\usepackage{graphicx}
\usepackage{multirow}
\usepackage{booktabs}
\usepackage{float}
\usepackage{setspace}
\usepackage{caption}
\usepackage{amsmath}
\usepackage{geometry}
\usepackage{lineno}
\usepackage{enumitem}

% bibliography/citations
\usepackage[super, comma, sort&compress]{natbib}

% hyperlinks should usually come after citation packages
\usepackage[colorlinks=true,linkcolor=black,citecolor=black,urlcolor=black]{hyperref}

\captionsetup[figure]{name=Figure,labelsep=period}
\captionsetup[table]{name=Table,labelsep=period}
\captionsetup{justification=centering}

\doublespacing
\setlength{\parindent}{0pt}

% Spacing Rules 

% Compact equation spacing
\setlength{\abovedisplayskip}{6pt}
\setlength{\belowdisplayskip}{4pt}
\setlength{\abovedisplayshortskip}{4pt}
\setlength{\belowdisplayshortskip}{4pt}

% Compact subsubsection spacing
\usepackage{titlesec}
\titlespacing*{\subsubsection}{0pt}{8pt}{4pt}




\geometry{
    left=2.54cm,
    right=2.54cm,
    top=2cm,
    bottom=2.54cm
}


\begin{document}

\section{Calculations}

\subsection{Dialysate Outlet Concentration, $C_{d,o}$}

Using the total urea mass balance,
\[
    Q_b(C_{b,i}-C_{b,o}) = Q_d(C_{d,o}-C_{d,i})
\]

Since fresh dialysate enters with negligible urea concentration, $C_{d,i}=0$. Solving for $C_{d,o}$:
\[
    C_{d,o} = \frac{Q_b}{Q_d}(C_{b,i}-C_{b,o})
    = \frac{1000}{2400}(2000-600)
\]

\[
    \boxed{C_{d,o}=583.3 \ \text{mg/L}}
\]

\subsection{Temperature-Corrected Urea Diffusivity, $D_{AB,T}$}

The diffusivity is corrected using:
\[
    D_{AB,T}
    = D_{AB,ref}
    \left(\frac{T}{T_{ref}}\right)
    \left(\frac{\mu_{ref}}{\mu_T}\right)
\]

The temperatures are converted to Kelvin:
\[
    T_{ref}=26.7+273.15=299.85 \ \text{K},
    \qquad
    T=37.4+273.15=310.55 \ \text{K}
\]

Substituting values:
\[
    D_{AB,T}
    = (1.44 \times 10^{-9})
    \left(\frac{310.55}{299.85}\right)
    \left(\frac{878.263 \times 10^{-6}}{699.037 \times 10^{-6}}\right)
\]

\[
    D_{AB,T}
    = (1.44 \times 10^{-9})(1.0357)(1.2564)
\]

\[
    \boxed{D_{AB,T}=1.87 \times 10^{-9} \ \text{m}^2/\text{s}}
\]

\subsection{Effective Diffusivity, $D_{Ae}$}

The effective diffusivity through the membrane is:
\[
    D_{Ae} = \varepsilon D_{AB,T}
    = (0.20)(1.87 \times 10^{-9})
\]

\[
    \boxed{D_{Ae}=3.74 \times 10^{-10} \ \text{m}^2/\text{s}}
\]

\subsection{Overall Mass Transfer Coefficient, $K_o$}

Assuming membrane resistance dominates:
\[
    K_o \approx \frac{D_{Ae}}{\delta_m}
    = \frac{3.74 \times 10^{-10}}{25 \times 10^{-6}}
\]

\[
    \boxed{K_o=1.50 \times 10^{-5} \ \text{m/s}}
\]

\subsection{Concentration Differences, $\Delta C_1$ and $\Delta C_2$}

For counter-current flow:
\[
    \Delta C_1 = C_{b,i}-C_{d,o},
    \qquad
    \Delta C_2 = C_{b,o}-C_{d,i}
\]

Since $C_{d,i}\approx 0$:
\[
    \Delta C_1 = 2000-583.3 = 1416.7 \ \text{mg/L}
\]

\[
    \Delta C_2 = 600-0 = 600 \ \text{mg/L}
\]

\[
    \boxed{\Delta C_1=1416.7 \ \text{mg/L}},
    \qquad
    \boxed{\Delta C_2=600 \ \text{mg/L}}
\]

\subsection{Log-Mean Concentration Difference, $\Delta C_{lm}$}

The log-mean concentration difference is:
\[
    \Delta C_{lm}
    =
    \frac{\Delta C_1-\Delta C_2}
    {\ln\left(\frac{\Delta C_1}{\Delta C_2}\right)}
\]

Using $\Delta C_1=1416.7 \ \text{mg/L}$ and $\Delta C_2=600 \ \text{mg/L}$:
\[
    \Delta C_{lm}
    =
    \frac{1416.7-600}
    {\ln\left(\frac{1416.7}{600}\right)}
    =
    \frac{816.7}{0.859}
\]

\[
    \boxed{\Delta C_{lm} \approx 951 \ \text{mg/L}}
\]

\subsection{Urea Removal Rate, $\dot{m}_A$}

The mass removal rate is:
\[
    \dot{m}_A = Q_b(C_{b,i}-C_{b,o})
    = (1.0)(2000-600)
    = 1400 \ \text{mg/min}
\]

\[
    \boxed{\dot{m}_A = 1400 \ \text{mg/min}}
\]

\subsection{Required Membrane Surface Area, $A_m$}

The required membrane area is:
\[
    A_m =
    \frac{Q_b(C_{b,i}-C_{b,o})}
    {K_o \Delta C_{lm}}
\]

Using $Q_b=1.0 \ \text{L/min}=1.667 \times 10^{-5} \ \text{m}^3/\text{s}$:
\[
    A_m
    =
    \frac{(1.667 \times 10^{-5})(1400)}
    {(1.50 \times 10^{-5})(951)}
    =
    \frac{2.334 \times 10^{-2}}
    {1.4265 \times 10^{-2}}
\]

\[
    \boxed{A_m \approx 1.64 \ \text{m}^2}
\]

\subsection{Number of Hollow Fibers, $N_f$}

The membrane area is related to hollow fiber geometry by:
\[
    A_m = N_f \pi d_f L
\]

Solving for $N_f$:
\[
    N_f = \frac{A_m}{\pi d_f L}
\]

Using $A_m=1.64 \ \text{m}^2$, $d_f=180 \ \mu\text{m}=1.80 \times 10^{-4} \ \text{m}$, and $L=0.20 \ \text{m}$:
\[
    N_f
    =
    \frac{1.64}
    {\pi(1.80 \times 10^{-4})(0.20)}
    =
    \frac{1.64}{1.131 \times 10^{-4}}
\]

\[
    \boxed{N_f \approx 1.45 \times 10^4 \ \text{fibers}}
\]

\subsection{Membrane Cost}

The total membrane cost is:
\[
    \text{Cost}
    =
    A_m \left(\frac{\$}{\text{m}^2}\right)
\]

Using $A_m=1.64 \ \text{m}^2$ and a membrane cost of $15 \ \$/\text{m}^2$:
\[
    \text{Cost} = 1.64(15)
\]

\[
    \boxed{\text{Cost} \approx \$24.60}
\]

\end{document}